\documentclass[10pt,a4paper]{article}
\usepackage[T1]{fontenc}
\usepackage[utf8]{inputenc}
\usepackage[margin=2cm]{geometry}
\usepackage{amsmath,amssymb}
\usepackage{longtable}
\usepackage{array}
\usepackage{booktabs}
\usepackage[colorlinks,linkcolor=blue!50!black,urlcolor=blue!50!black]{hyperref}

\newcolumntype{S}{>{\raggedright\arraybackslash$}p{0.25\textwidth}<{$}}
\newcolumntype{M}{>{\raggedright\arraybackslash}p{0.51\textwidth}}
\newcolumntype{L}{>{\raggedright\arraybackslash}p{0.16\textwidth}}

\newcommand{\symtable}[1]{%
  \begin{longtable}{SML}
  \toprule
  \multicolumn{1}{l}{\textbf{Symbol}} & \textbf{Meaning} & \textbf{Where} \\
  \midrule
  \endfirsthead
  \toprule
  \multicolumn{1}{l}{\textbf{Symbol}} & \textbf{Meaning} & \textbf{Where} \\
  \midrule
  \endhead
  \bottomrule
  \endfoot
  #1
  \end{longtable}%
}

\title{The Matrix Cookbook (Petersen \& Pedersen)\\[2pt]
       \large Summary of Major Symbols and Concepts, by Chapter}
\author{Summary of \texttt{matrix-cookbook.pdf}, version November 15, 2012 (72 pp.)}
\date{}

\begin{document}
\maketitle
\tableofcontents
\bigskip

\section*{What this is --- and what it is not}
\addcontentsline{toc}{section}{What this is --- and what it is not}

The Cookbook is a \emph{reference}, not a text.  It is a collection of identities,
approximations, inequalities and relations about matrices, numbered
consecutively (1--\,$\sim$550) across the whole document so that any single result
can be cited by equation number.  There are almost no proofs --- Appendix~B holds
the handful that exist --- and the authors state plainly that the material was
``collected, borrowed and copied'' and very likely contains errors.

Practical consequences of that design:

\begin{itemize}
\item \textbf{Cite by equation number, not by page.}  Equation numbers are stable
      across the document; section numbering is not especially memorable.
\item \textbf{Check the layout convention before transplanting a derivative.}  The
      Cookbook's $\partial/\partial\mathbf{X}$ results are laid out so that
      $\partial\det(\mathbf{X})/\partial\mathbf{X}=\det(\mathbf{X})(\mathbf{X}^{-1})^{T}$
      (Eq.~49) --- i.e.\ the result has the \emph{same} shape as $\mathbf{X}$.  This
      is the ``gradient'' convention, and it is transposed relative to the numerator
      layout used in, e.g., Dey's note.
\item \textbf{Verify anything load-bearing.}  The authors solicit corrections; the
      document is explicitly a best-effort compilation.
\end{itemize}

%%%%%%%%%%%%%%%%%%%%%%%%%%%%%%%%%%%%%%%%%%%%%%%%%%%%%%%%%%%%%%%%%%%%%%
\section{Notation and nomenclature (front matter, p.~5)}

The Cookbook prints its own notation table.  Reproduced and grouped here.

\subsection*{Objects and indexing}
\symtable{
\mathbf{A}                      & a matrix                                                       & p.~5 \\
a,\;\mathbf{a}                  & a scalar; a (column) vector                                    & p.~5 \\
\mathbf{A}_{ij},\;\mathbf{A}_{i},\;\mathbf{A}^{ij} & a matrix indexed for some purpose (not an entry) & p.~5 \\
\mathbf{A}^{n}                  & the $n$-th power of a square matrix (or an indexed matrix)     & p.~5 \\
(\mathbf{A})_{ij},\;A_{ij}      & the $(i,j)$-th \emph{entry} of $\mathbf{A}$                     & p.~5 \\
{[\,\mathbf{A}\,]}_{ij}         & the $ij$-\textbf{submatrix}: $\mathbf{A}$ with row $i$ and column $j$ deleted & p.~5 \\
a_{i}                           & the $i$-th element of the vector $\mathbf{a}$                  & p.~5 \\
}

\subsection*{Derived matrices}
\symtable{
\mathbf{A}^{-1}                 & inverse                                                        & \S3 \\
\mathbf{A}^{+}                  & the \textbf{pseudo (Moore--Penrose) inverse}                   & \S3.6 \\
\mathbf{A}^{1/2}                & matrix square root (if unique) --- \emph{not} elementwise       & p.~5 \\
\mathbf{A}^{T}                  & transpose                                                      & p.~5 \\
\mathbf{A}^{-T}                 & $(\mathbf{A}^{-1})^{T}=(\mathbf{A}^{T})^{-1}$                  & p.~5 \\
\mathbf{A}^{*}                  & complex conjugate (entrywise)                                  & \S4 \\
\mathbf{A}^{H}                  & conjugate transpose (Hermitian adjoint)                        & \S4 \\
}

\subsection*{Scalar-valued operators}
\symtable{
\det(\mathbf{A})                & determinant                                                    & \S1.2 \\
\mathrm{Tr}(\mathbf{A})         & trace                                                          & \S1.1 \\
\mathrm{diag}(\mathbf{A})       & the diagonal matrix with $(\mathrm{diag}(\mathbf{A}))_{ij}=\delta_{ij}A_{ij}$ & p.~5 \\
\mathrm{eig}(\mathbf{A})        & the eigenvalues of $\mathbf{A}$                                & \S5.2 \\
\mathrm{vec}(\mathbf{A})        & the vector form of $\mathbf{A}$ (columns stacked)              & \S10.2.2 \\
\lVert\mathbf{A}\rVert          & a matrix norm; a subscript names which                         & \S10.4 \\
\sup                            & supremum of a set                                              & p.~5 \\
\Re z,\;\Re\mathbf{Z}           & real part of a scalar / vector / matrix                        & \S4 \\
\Im z,\;\Im\mathbf{Z}           & imaginary part                                                 & \S4 \\
}

\subsection*{Products and named matrices}
\symtable{
\mathbf{A}\circ\mathbf{B}       & \textbf{Hadamard} (elementwise) product                        & \S10.2 \\
\mathbf{A}\otimes\mathbf{B}     & \textbf{Kronecker} product                                     & \S10.2 \\
\mathbf{0}                      & the null matrix (zero in all entries)                          & p.~5 \\
\mathbf{I}                      & the identity matrix                                            & p.~5 \\
\mathbf{J}^{ij}                 & the \textbf{single-entry matrix}: $1$ at $(i,j)$, zero elsewhere & \S9.7 \\
\mathbf{\Sigma}                 & a positive definite matrix (typically a covariance)            & \S6--8 \\
\mathbf{\Lambda}                & a diagonal matrix (typically of eigenvalues)                   & \S5.2 \\
\mathbf{S}^{ij}                 & the \textbf{structure matrix}, $\partial\mathbf{A}/\partial A_{ij}$; equals $\mathbf{J}^{ij}$ when $\mathbf{A}$ is unstructured & \S2.8 \\
}

%%%%%%%%%%%%%%%%%%%%%%%%%%%%%%%%%%%%%%%%%%%%%%%%%%%%%%%%%%%%%%%%%%%%%%
\section{Chapter map with headline results}

\subsection{1. Basics (p.~6)}
Trace (\S1.1), determinant (\S1.2), and the fully written-out $2\times2$ case
(\S1.3).  The standard facts: $\mathrm{Tr}(\mathbf{AB})=\mathrm{Tr}(\mathbf{BA})$,
$\mathrm{Tr}(\mathbf{A})=\sum_{i}\lambda_{i}$,
$\det(\mathbf{AB})=\det(\mathbf{A})\det(\mathbf{B})$,
$\det(\mathbf{A})=\prod_{i}\lambda_{i}$.

\subsection{2. Derivatives (p.~8) --- the most-used chapter}
\begin{description}
\item[\S2.1 Determinant.]
      $\dfrac{\partial\det(\mathbf{X})}{\partial\mathbf{X}}=\det(\mathbf{X})(\mathbf{X}^{-1})^{T}$
      (Eq.~49); $\dfrac{\partial\ln\lvert\det(\mathbf{X})\rvert}{\partial\mathbf{X}}=(\mathbf{X}^{-1})^{T}$
      (Eq.~57); square forms $\det(\mathbf{X}^{T}\mathbf{A}\mathbf{X})$ split into
      three cases according to whether $\mathbf{X}$ is square and $\mathbf{A}$
      symmetric (Eqs.~52--54).
\item[\S2.2 Inverse.]  The one identity everything else rests on,
      $\dfrac{\partial\mathbf{Y}^{-1}}{\partial x}=-\mathbf{Y}^{-1}\dfrac{\partial\mathbf{Y}}{\partial x}\mathbf{Y}^{-1}$
      (Eq.~59).
\item[\S2.3 Eigenvalues.]  $\partial\lambda_{i}=\mathbf{v}_{i}^{T}\partial(\mathbf{A})\mathbf{v}_{i}$
      and $\partial\mathbf{v}_{i}=(\lambda_{i}\mathbf{I}-\mathbf{A})^{+}\partial(\mathbf{A})\mathbf{v}_{i}$
      for real symmetric $\mathbf{A}$ with distinct eigenvalues (Eqs.~67--68).
\item[\S2.4 Matrices, vectors, scalar forms.]  The bulk lookup table --- first-order,
      second-order and higher-order forms.
\item[\S2.5 Traces.]  A long list; the workhorses are
      $\partial\,\mathrm{Tr}(\mathbf{X}^{T}\mathbf{X})/\partial\mathbf{X}=2\mathbf{X}$
      (Eq.~115) and
      $\partial\,\mathrm{Tr}(\mathbf{AXBX}^{T}\mathbf{C})/\partial\mathbf{X}=\mathbf{A}^{T}\mathbf{C}^{T}\mathbf{X}\mathbf{B}^{T}+\mathbf{CAXB}$
      (Eq.~118).
\item[\S2.6--2.7 Norms.]  Vector norms; matrix norms, e.g.\
      $\partial\lVert\mathbf{X}\rVert_{F}^{2}/\partial\mathbf{X}=2\mathbf{X}$ (Eq.~132).
\item[\S2.8 Structured matrices --- read this before using \S2.4--2.5.]
      If $\mathbf{A}$ is symmetric, Toeplitz, etc., \emph{the preceding sections do
      not apply}.  The general rule is
      $\dfrac{df}{dA_{ij}}=\mathrm{Tr}\!\left[\left(\dfrac{\partial f}{\partial\mathbf{A}}\right)^{T}\dfrac{\partial\mathbf{A}}{\partial A_{ij}}\right]$
      (Eq.~133), with $\partial\mathbf{A}/\partial A_{ij}=\mathbf{S}^{ij}$ the
      \textbf{structure matrix} (Eq.~134).  For symmetric $\mathbf{A}$,
      $\mathbf{S}^{ij}=\mathbf{J}^{ij}+\mathbf{J}^{ji}-\mathbf{J}^{ij}\mathbf{J}^{ij}$,
      giving
      $\dfrac{df}{d\mathbf{A}}=\dfrac{\partial f}{\partial\mathbf{A}}+\left(\dfrac{\partial f}{\partial\mathbf{A}}\right)^{T}-\mathrm{diag}\!\left(\dfrac{\partial f}{\partial\mathbf{A}}\right)$
      (Eq.~138).  \S2.8.1 gives the matrix chain rule
      $\dfrac{\partial g(\mathbf{U})}{\partial X_{ij}}=\mathrm{Tr}\!\left[\left(\dfrac{\partial g(\mathbf{U})}{\partial\mathbf{U}}\right)^{T}\dfrac{\partial\mathbf{U}}{\partial X_{ij}}\right]$
      (Eq.~137).
\end{description}

\subsection{3. Inverses (p.~17)}
Basic facts; \textbf{exact relations} (\S3.2) --- the family of matrix-inversion
lemmas that make large inverses tractable:
\begin{itemize}
\item \textbf{Woodbury} (Eqs.~156--157):
      $(\mathbf{A}+\mathbf{UBV})^{-1}=\mathbf{A}^{-1}-\mathbf{A}^{-1}\mathbf{U}(\mathbf{B}^{-1}+\mathbf{VA}^{-1}\mathbf{U})^{-1}\mathbf{VA}^{-1}$.
\item \textbf{Kailath variant} (Eq.~159):
      $(\mathbf{A}+\mathbf{BC})^{-1}=\mathbf{A}^{-1}-\mathbf{A}^{-1}\mathbf{B}(\mathbf{I}+\mathbf{CA}^{-1}\mathbf{B})^{-1}\mathbf{CA}^{-1}$.
\item \textbf{Sherman--Morrison} (Eq.~160), the rank-one case:
      $(\mathbf{A}+\mathbf{bc}^{T})^{-1}=\mathbf{A}^{-1}-\dfrac{\mathbf{A}^{-1}\mathbf{bc}^{T}\mathbf{A}^{-1}}{1+\mathbf{c}^{T}\mathbf{A}^{-1}\mathbf{b}}$.
\item \textbf{The Searle set} (Eqs.~161--167), including
      $(\mathbf{I}+\mathbf{AB})^{-1}=\mathbf{I}-\mathbf{A}(\mathbf{I}+\mathbf{BA})^{-1}\mathbf{B}$.
\end{itemize}
Then implications on inverses (\S3.3), approximations (\S3.4), generalized inverses
(\S3.5), and the \textbf{pseudo inverse} $\mathbf{A}^{+}$ (\S3.6) with rank-one
update formulas.

\subsection{4. Complex matrices (p.~24)}
Complex derivatives (\S4.1) --- note the modified chain rule for \emph{non-analytic}
$f$, which reduces to the ordinary one in the analytic case; higher-order and
non-linear derivatives (\S4.2); inverse of a complex sum (\S4.3).

\subsection{5. Solutions and decompositions (p.~28)}
Solutions to linear equations (\S5.1); eigenvalues and eigenvectors (\S5.2);
\textbf{SVD} (\S5.3); triangular, \textbf{LU}, \textbf{LDM}, \textbf{LDL}
decompositions (\S5.4--5.7).

\subsection{6--8. Statistics, distributions, Gaussians (pp.~34--45)}
Definition of moments (\S6.1); expectation of linear combinations (\S6.2); weighted
scalar variables (\S6.3).  Multivariate distributions (\S7): Cauchy, Dirichlet,
Normal, Normal-Inverse Gamma, Gaussian, Multinomial, Student's $t$, Wishart and
inverse Wishart.  Chapter 8 is the Gaussian workhorse: products, marginals and
conditionals, linear combinations, rearranging means, moments, and mixtures.  These
three chapters are why the Cookbook is a fixture in machine-learning work.

\subsection{9. Special matrices (p.~46)}
\begin{itemize}
\item \textbf{Block matrices} (\S9.1), including the \textbf{Schur complement}
      (\S9.1.5) and the block-inverse and block-determinant formulas built from it.
      The Cookbook notes (\S8) that Gaussian conditional covariances \emph{are} Schur
      complements --- the link between chapters 8 and 9.
\item Discrete Fourier transform matrix (\S9.2); Hermitian and skew-Hermitian
      (\S9.3); idempotent (\S9.4); orthogonal (\S9.5); positive (semi-)definite
      (\S9.6, with the Hadamard inequality); the \textbf{single-entry matrix}
      $\mathbf{J}^{ij}$ (\S9.7, the atom of all the structure-matrix results);
      symmetric and skew-symmetric (\S9.8); Toeplitz (\S9.9); transition matrices
      (\S9.10); units, permutation and shift (\S9.11); Vandermonde (\S9.12).
\end{itemize}

\subsection{10. Functions and operators (p.~58)}
Functions and series of matrices (\S10.1); \textbf{Kronecker and vec} (\S10.2), whose
central identity is
\[
\mathrm{vec}(\mathbf{AXB})=(\mathbf{B}^{T}\otimes\mathbf{A})\,\mathrm{vec}(\mathbf{X})
\quad(\text{Eq.}~520),
\qquad
\mathrm{Tr}(\mathbf{A}^{T}\mathbf{B})=\mathrm{vec}(\mathbf{A})^{T}\mathrm{vec}(\mathbf{B})
\quad(\text{Eq.}~521);
\]
vector norms (\S10.3); matrix norms (\S10.4); rank (\S10.5); integrals involving Dirac
deltas (\S10.6); miscellaneous (\S10.7).

\subsection{Appendices}
\textbf{A} --- one-dimensional results (Gaussian; 1-D mixture of Gaussians), as a
sanity check on the matrix versions.  \textbf{B} --- the few proofs and details.

%%%%%%%%%%%%%%%%%%%%%%%%%%%%%%%%%%%%%%%%%%%%%%%%%%%%%%%%%%%%%%%%%%%%%%
\section{Notation traps worth flagging}

\begin{itemize}
\item \textbf{$\mathbf{A}_{ij}$ is ambiguous by design.}  The notation table lists
      $\mathbf{A}_{ij}$ both as ``a matrix indexed for some purpose'' and as the
      $(i,j)$-th entry.  Context decides; bold weight is the only cue.
\item \textbf{$[\mathbf{A}]_{ij}$ is a submatrix, not an entry} --- it is $\mathbf{A}$
      with row $i$ and column $j$ \emph{deleted}.  Adjacent to $(\mathbf{A})_{ij}$,
      which \emph{is} the entry.
\item \textbf{$\mathbf{A}^{1/2}$ is a matrix square root}, not the elementwise square
      root; and it need not be unique.
\item \textbf{$\mathbf{A}^{n}$ collides with the ``indexed matrix'' convention.}  A
      superscript may be a power or a label.
\item \textbf{Layout.}  Cookbook derivative results are shaped like the differentiated
      matrix (gradient convention).  Combining them with a numerator-layout source
      without transposing is the standard way to get a shape-wrong answer.
\item \textbf{\S2.4--2.7 assume no structure.}  For symmetric, Toeplitz or otherwise
      constrained matrices you must go through \S2.8; using the unstructured result
      double-counts the off-diagonal entries.
\item \textbf{Errors are expected.}  The authors say so explicitly and give a
      correction address.  Treat any single equation as a strong hint, not a proof.
\end{itemize}

\end{document}
